% Source: https://artofproblemsolving.com/community/c3h1671246_one_number_theory_problem
% Stable identity: aops-topic-1671246-q01
\textbf{Complete the square.} Multiply by $17$:
\[
17\left(17y^2+22y+1\right)=(17y+11)^2-121+17=(17y+11)^2-104 .
\]
So if $17y^2+22y+1=k^2$ and $X=17y+11$, then
\[
X^2-17k^2=104. \tag{$*$}
\]
Conversely, every integer solution $(X,k)$ of $(*)$ with $X\equiv 11\pmod{17}$ gives the integer $y=\dfrac{X-11}{17}$ with $17y^2+22y+1=\dfrac{X^2-104}{17}=k^2$.

\textbf{A map that produces new solutions.} For integers $X,k$ put
\[
X'=2177X+8976k,\qquad k'=528X+2177k .
\]
Expanding,
\[
X'^2-17k'^2=\left(2177^2-17\cdot528^2\right)X^2+2\left(2177\cdot8976-17\cdot528\cdot2177\right)Xk+\left(8976^2-17\cdot2177^2\right)k^2 .
\]
Here $2177^2-17\cdot528^2=4739329-4739328=1$, the middle bracket is $0$ because $8976=17\cdot528$, and $8976^2-17\cdot2177^2=17\left(17\cdot528^2-2177^2\right)=-17$. Hence $X'^2-17k'^2=X^2-17k^2$, so $(X',k')$ is again a solution of $(*)$. (The numbers come from $(33+8\sqrt{17})^2=2177+528\sqrt{17}$, and $33^2-17\cdot8^2=1$.)

Moreover $2177=17\cdot128+1$ and $17\mid 8976$, so $X'\equiv X\pmod{17}$: the condition $X\equiv11\pmod{17}$ is kept. Finally, if $X,k>0$ then $X'>X>0$ and $k'>0$.

\textbf{Infinitely many $y$.} Start from $y=0$, that is $(X,k)=(11,1)$. Applying the map again and again gives solutions with $X\equiv 11\pmod{17}$ and strictly increasing $X$, hence infinitely many different integers $y=\dfrac{X-11}{17}$. The first step gives $X'=2177\cdot11+8976=32923$ and $k'=528\cdot11+2177=7985$, so
\[
y=\dfrac{32923-11}{17}=1936,\qquad 17\cdot1936^2+22\cdot1936+1=7985^2 .
\]
So $17y^2+22y+1$ is a perfect square for infinitely many integers $y$. $\blacksquare$

(The values with $|y|<10^6$ are $y=-127790,\,-58110,\,-5586,\,-306,\,-30,\,-14,\,-2,\,0,\,4,\,84,\,880,\,1936,\,20148,\,368548$.)
